\input{preamble.tex}

\title{Отчет по лабораторной работе № 2. \\ Предварительная настройка оборудования Cisco}
\author{Сущенко Алина \\ НПИбд-01-23}
\date{2026}

\begin{document}

\maketitle
\newpage

\tableofcontents

\newpage
\section{Цель работы}
Получить основные навыки по начальному конфигурированию оборудования Cisco.

\newpage
\section{Выполнение работы}
\item В логической рабочей области \texttt{Packet Tracer} разместим коммутатор, маршрутизатор и 2 оконечных устройства типа PC( PC0 и PC1), соединим один PC с маршрутизатором, другой PC — с коммутатором, как показано на схеме. (Рис. \ref{01}):
\begin{center}
    \centering
    \includegraphics[width=\textwidth]{../images/01.png}
    \captionof{figure}{Cхема подключения оборудования для проведения его предварительной настройки.}
    \label{01}
\end{center}

\subsection{Конфигурация маршрутизатора}
\begin{enumerate}
\item Произведём настройку маршрутизитора. Присвоим устройству имя \texttt{msk-ansusenko-gw-1},настройим интерфейс FastEthernet 0/0, удалённый доступ и SSH. (Рис. \ref{02}):
\begin{center}
    \centering
    \includegraphics[width=\textwidth]{../images/02.png}
    \captionof{figure}{Настройка маршрутизатора.}
    \label{02}
\end{center}

\subsection{Конфигурация коммутатора}
\item Далее настроим коммутатор. Так же зададим ему имя \texttt{msk-ansusenko-sw-1}, создим и настроим VLAN 2, порт доступа и шлюз по умолчанию. Затем повторим то же, что делали с маршрутизитором.(Рис. \ref{03}):

\begin{center}
    \centering
    \includegraphics[width=\textwidth]{../images/03.png}
    \captionof{figure}{Настройка коммутатора.}
    \label{03}
\end{center}

\item Далее проверим, чтобы наши настройки применилсь. После этого проверим IP наших конечных устройств. (Рис. \ref{04}, \ref{05}):
\begin{center}
    \centering
    \includegraphics[width=\textwidth]{../images/04.png}
    \captionof{figure}{Настройка РС0.}
    \label{04}
\end{center}

\begin{center}
    \centering
    \includegraphics[width=\textwidth]{../images/05.png}
    \captionof{figure}{Настройка РС1.}
    \label{05}
\end{center}

\item Теперь проверим связь при помощи наших РС. На РС0 выполним команды: (Рис. \ref{06}):
  \begin{minted}{bash}
    ping 192.168.2.10
    telnet 192.168.2.1
    ssh -l admin 192.168.2.1
  \end{minted}

\begin{center}
    \centering
    \includegraphics[width=\textwidth]{../images/06.png}
    \captionof{figure}{Проверка подключения и настроек при помощи команд на РС0.}
    \label{06}
\end{center}

\item Затем то же самое, но для РС1 (Рис. \ref{07}):
\begin{minted}{bash}
    ping 192.168.1.254
    telnet 192.168.1.254
    ssh -l admin 192.168.1.254
  \end{minted}
\begin{center}
    \centering
    \includegraphics[width=\textwidth]{../images/07.png}
    \captionof{figure}{Проверка подключения и настроек при помощи команд на РС1.}
    \label{07}
\end{center}

\end{enumerate}

\section{Выводы}
В результате выполнения работы мы успешно воссоздали схему и настроили её. Так же получили основные навыки по начальному конфигурированию оборудования Cisco.

\end{document}
